\documentclass[12pt,a4paper]{ctexart}
\usepackage[margin=2.5cm]{geometry}
\usepackage{graphicx}
\usepackage{booktabs}
\usepackage{longtable}
\usepackage{hyperref}
\usepackage{xcolor}
\usepackage{listings}
\usepackage{setspace}
\onehalfspacing

\hypersetup{
  colorlinks=true,
  linkcolor=blue,
  urlcolor=blue,
  citecolor=blue
}

\title{\textbf{smallgameagent 实验报告} \\ \large 基于 LLM/VLM 与规则的小游戏 Agent}
\author{smallgameagent 项目组}
\date{\today}

\begin{document}
\maketitle
\tableofcontents
\newpage
\section{smallgameagent 实验报告（终稿）}

> 2026-07-17。配套：\textbackslash\{\}texttt\{EXPERIMENT\textbackslash\{\}\_PLAN.md\}（方案）、\textbackslash\{\}texttt\{EXPERIMENT\textbackslash\{\}\_RESULTS.md\}（过程数据）。\\
> 本文面向合作同学与老师，按四个问题组织结论，附三条决策路线的效率-质量画像与训练准备状态。\\

\subsection{0. 一页结论}

\begin{itemize}
  \item \textbackslash\{\}textbf\{空间一致性\}：根因不是「记不住场景」而是\textbackslash\{\}textbf\{交互方式随场景切换失效未被检测\}。实证：
\end{itemize}
00461 场景切换后 \textbackslash\{\}texttt\{LosePanel\} 激活阻断摇杆输入，规则 agent 继续发移动指令瘫痪 175+ 步。\\
解法是\textbackslash\{\}textbf\{版本化世界模型 + 失效检测与恢复\}（stale 标记、失败面板自动点击），\\
A/B 显示活动率 0.54→0.92、墙钟 263→154s（−42\%）。\\
\begin{itemize}
  \item \textbackslash\{\}textbf\{时间一致性\}：交付\textbackslash\{\}textbf\{阶段契约机制\}（三层时间对齐 + settle 复查 + 受保护前缀 hash +
\end{itemize}
失败分级 rollback/compensation/stop+replan），29 单测含真实温度案例复刻（131/190 步偏差）。\\
\begin{itemize}
  \item \textbackslash\{\}textbf\{策略优化\}：kimi-k2.7-code 的经济决策实验显示\textbackslash\{\}textbf\{解析与 token 预算才是第一瓶颈\}
\end{itemize}
（4K max\_tokens 下 58\% 调用截断无答案）；修正后朴素与引导式在饱和场景打平（0.93），\\
决定性场景（batch 占优 3-4×）结果见 §3。\\
\begin{itemize}
  \item \textbackslash\{\}textbf\{效率\}：动态探针预算模块（五触发器、L0/L1/L2 分层、L2 预算上限 20\%）+ 动态日志分级，
\end{itemize}
27 单测，模拟装饰性场景节省观测成本 \textasciitilde{}82\%。\\
\begin{itemize}
  \item \textbackslash\{\}textbf\{VLM 本地推理\}：Intel 核显（WSL2 Mesa Dozen/Vulkan）文本基准实测：
\end{itemize}
\textbackslash\{\}textbf\{Qwen3.5-4B 2.75 tok/s、Qwen3.5-9B 0.72 tok/s、gemma-4-E4B 0.67 tok/s\}（Q4\_K\_M，-ngl 99）。\\
视觉编码必须 \textbackslash\{\}texttt\{--no-mmproj-offload\} 回 CPU，否则挂死或输出空白；4B 视觉 struct 4 帧评测\\
\textbackslash\{\}textbf\{解析率 0/4\}（超时/连接抖动），平均 534 s/帧，\textbackslash\{\}textbf\{当前仅适合离线标注，不适合在线逐步控制\}。\\
\begin{itemize}
  \item \textbackslash\{\}textbf\{训练\}：processed-runs 已转为 7 任务 15,083 样本（14 单测+加载验证），
\end{itemize}
训练脚本就绪，等 ssh5090 可访问后开 QLoRA。\\

\subsection{1. 环境打通（新信息）}

\begin{table}[htbp]
\centering
\small
\begin{tabular}{|l|l|}
\hline
项 & 结论 \\
\hline
云端 LLM & \textbackslash\{\}texttt\{https://opencode.ai/zen/go/v1\}，kimi-k2.7-code 文本决策 2.4–6.3s/次（旧 deepseek 8–15s），多模态可用；\textbackslash\{\}textbf\{坑：显式 temperature 参数触发 400\}，已在客户端按模型名省略 \\
\hline
本地 VLM & LM Studio 捆绑 llama.cpp Vulkan 后端（WSL2 Mesa Dozen→Intel Graphics）；lms CLI 因版本错位无法认证，改直跑 \textbackslash\{\}texttt\{llama-server\}（模型放 \textbackslash\{\}texttt\{\textbackslash\{\}textasciitilde\{\}/.lmstudio/models/\}，与 LM Studio 应用互通） \\
\hline
游戏基线 & rule 模式 SSD\_00461 塔防 300 步基线 composite 0.425（verifiers 风格 rubric） \\
\hline
\end{tabular}
\end{table}

\subsection{2. 空间一致性（问题①）}

\textbackslash\{\}textbf\{实证根因\}（三级递进，全部有 run 数据）：\\
1. 初判「场景切换后撞新障碍卡死」→ 障碍学习+势场避让（P1）上线后仍失败；\\
2. 再判「英雄死亡」→ 诊断跑血量全程 100，排除；\\
3. \textbackslash\{\}textbf\{定案：failCount 翻转 → LosePanel 激活 → 输入被面板阻断，agent 不点继续 → 永久瘫痪\}。\\
这正是同学说的「交互方式改变」类失效（摇杆交互被模态面板替换）。\\

\textbackslash\{\}textbf\{修复\}：\\
\begin{itemize}
  \item \textbackslash\{\}texttt\{VersionedWorldModel\}（\textbackslash\{\}texttt\{src/agent/world\textbackslash\{\}\_model.py\}）：entity\_version/scene\_epoch/capability\_epoch，
\end{itemize}
观测写入比对、stale 标记、局部重规划，12 单测（含 autoFishing 12 次翻转复刻）。\\
\begin{itemize}
  \item 探针新增 \textbackslash\{\}texttt\{findPanelButtons\}：按\textbackslash\{\}textbf\{节点名\}匹配按钮（生产构建组件名被 minify 成单字母，
\end{itemize}
按类名匹配必然落空），返回设计分辨率坐标；Python 侧 retry 优选、避开 download 广告陷阱，\\
设计坐标→CSS 像素映射（720×1560 左下原点 → 375×812 顶左）。7 单测。\\
\begin{itemize}
  \item RuleEngine 障碍学习（势场避让+定向逃逸，15 单测）——对几何卡死有效，但非本游戏 binding。
\end{itemize}

\textbackslash\{\}textbf\{A/B（00461，300 步，rubric）\}：\\

\begin{table}[htbp]
\centering
\small
\begin{tabular}{|l|l|l|l|l|l|}
\hline
run & composite & activity & progress & stall 步 & 墙钟 \\
\hline
基线 & 0.425 & 0.535 & 0.667 & 139 & 262.9s \\
\hline
+世界模型/障碍 & 0.351 & 0.408 & 0.583 & 177 & 320.7s \\
\hline
\textbackslash\{\}textbf\{+失败面板处理\} & \textbackslash\{\}textbf\{0.473\} & \textbackslash\{\}textbf\{0.920\} & \textbackslash\{\}textbf\{0.750\} & \textbackslash\{\}textbf\{24\} & \textbackslash\{\}textbf\{153.5s\} \\
\hline
\end{tabular}
\end{table}

\subsection{3. 策略优化（问题③）}

实验：12 经济场景（gold/双升级项/收入/预算），最优解由穷举模拟器给出，kimi-k2.7-code\\
朴素 prompt vs 机制引导 prompt（先抽象经济循环+往返成本再决策）。\\

\begin{itemize}
  \item v1（max\_tokens 4096）：\textbackslash\{\}textbf\{14/24 次调用无答案\}——思考型模型把 token 预算耗尽于推理，
\end{itemize}
JSON 被截断。教训：思考模型必须留足输出预算（≥16K）并约束「末行只放 JSON」。\\
\begin{itemize}
  \item v2（16K tokens）：解析失败归零；朴素 vs 引导 0.929 vs 0.932（场景饱和，11/12 greedy 最优，无法区分）。
  \item v3（8 个 batch 占优 3-4× + 4 对照，模拟器穷举真值）：
\end{itemize}
\textbackslash\{\}textbf\{最优命中率 朴素 41.7\textbackslash\{\}\% vs 引导 66.7\textbackslash\{\}\%\}；按得分比率（所选策略得分/最优得分）：\\
**朴素 batch=0.711 / control=1.000（overall 0.808）；引导 batch=1.000 / control=0.958\\
（overall 0.986）**。\\
结论：机制引导把 batch 经济决策从 71\% 提升到 100\% 最优值，代价仅 greedy 场景 −4\%——\\
定量证实了「加以引导后探索最优解能力强」的判断。引导的关键是 prompt 里显式给出\\
「经济循环 + 往返/机会成本」的抽象，而不是让模型自己发现。\\

\subsection{4. 效率（问题④）}

\begin{itemize}
  \item \textbackslash\{\}texttt\{src/agent/probe\textbackslash\{\}\_budget.py\}：L0 状态（73ms）/L1 组件（150ms）/L2 截图（400ms）三级，
\end{itemize}
五触发器（phase\_change/low\_confidence/action\_no\_effect/collision\_hint/semantic\_flip），\\
冷却防抖 + L2 窗口占比上限（默认 20\%）+ 高优先挤占；AdaptiveLogger 环形内存 DEBUG +\\
触发窗口落盘。27 单测；50 步无变化场景模拟节省 81.75\%。\\
\begin{itemize}
  \item 失败面板自动恢复即「策略级回退」的首个实例（行为级：障碍逃逸；状态级：stale 重规划；
\end{itemize}
策略级：面板 dismiss）。\\

\subsection{5. 时间一致性（问题②）}

\textbackslash\{\}texttt\{src/agent/phase\textbackslash\{\}\_contract.py\}（29 单测）：TimestampedValue 三层时间戳（event/observed/settled）、\\
跨字段一致性校验器（复刻 190 步 131 违例）、PhaseGate 状态机（settle 复查防完成假阳性、\\
守卫违反 VIOLATED、超时 TIMEOUT）、受保护前缀 hash（篡改即 PrefixViolation）、\\
失败按副作用分级。与游戏 loop 的集成（契约守卫替换硬编码 fail 检测）是下一步。\\

\subsection{6. 本地 VLM 画像（P5）}

\subsubsection{6.1 文本生成基准（Q4\_K\_M，-ngl 99，prompt="1, 2, 3,"，max\_tokens=60）}

\begin{table}[htbp]
\centering
\small
\begin{tabular}{|l|l|l|l|}
\hline
模型 & server tok/s & wall 60 tokens & 备注 \\
\hline
Qwen3.5-4B & \textbackslash\{\}textbf\{2.75\} & 24.6 s & 可用 \\
\hline
Qwen3.5-9B & 0.72 & 90.3 s & 慢但可跑 \\
\hline
gemma-4-E4B-it & 0.67 & 37.5 s & 输出格式偏闲聊，需强约束 \\
\hline
\end{tabular}
\end{table}

\subsubsection{6.2 视觉 struct 离线评测}

\begin{itemize}
  \item 4 帧真实截图（\textbackslash\{\}texttt\{SSD\textbackslash\{\}\_00496P01\} step 2/11，\textbackslash\{\}texttt\{SSD\textbackslash\{\}\_00219P01\} step 1/10），
\end{itemize}
Qwen3.5-4B Q4\_K\_M + \textbackslash\{\}texttt\{--no-mmproj-offload\}：\\
\begin{itemize}
  \item \textbackslash\{\}textbf\{解析率 0/4\}，目标准确率 0/4，平均延迟 534 s/帧；
  \item 失败原因：单帧生成超时（>600 s）与 LM Studio 客户端连接抖动；
  \item 真值全部 \textbackslash\{\}texttt\{has\textbackslash\{\}\_target=True\}，模型即使输出也未能在时限内返回可解析 JSON。
  \item 单独用 \textbackslash\{\}texttt\{max\textbackslash\{\}\_tokens=256\} 重跑一帧 step 2 的测试见
\end{itemize}
\textbackslash\{\}texttt\{experiment\textbackslash\{\}\_vlm\textbackslash\{\}\_one\textbackslash\{\}\_frame\textbackslash\{\}\_4b.json\}（\textbackslash\{\}textbf\{78.7 s 后连接错误，仍未解析成功\}）。\\

\subsubsection{6.3 结论}

\begin{itemize}
  \item 本地 4B/9B/E4B 在 Intel 核显上\textbackslash\{\}textbf\{仅适合离线数据标注\}；在线逐步控制应继续使用
\end{itemize}
云端 mimo-v2.5（约 9.5 s/帧）或更强的本地 GPU。\\
\begin{itemize}
  \item 视觉管线瓶颈在 \textbackslash\{\}textbf\{CPU mmproj 编码\}（\textasciitilde{}30–90 s）和 \textbackslash\{\}textbf\{低 tok/s 生成\}，二者叠加后
\end{itemize}
单帧常超 5 min；QLoRA 微调前应先固定 \textbackslash\{\}texttt\{--no-mmproj-offload\} 与足够长的 HTTP 超时。\\

\subsection{7. 训练准备（P7，等 ssh5090）}

\begin{itemize}
  \item \textbackslash\{\}texttt\{vlm-training-data-processed-runs/\}：22 游戏 3054 步 → \textbackslash\{\}textbf\{15,083 样本 / 7 任务\}
\end{itemize}
（next\_probe\_action 2645 / probe\_action\_effect 2645 / field\_grounding 2645 /\\
information\_gain 3054 / pulse\_response 1435 / progression 2645 / failure\_recovery 14），\\
VLMColdStartDataset 加载验证通过，14 单测。\\
\begin{itemize}
  \item failure\_recovery 仅 14 条（成功 run 中卡死窗口少）——建议用 verifiers 环境跑
\end{itemize}
对抗性失败 run 补数据（rule 变体/随机扰动注入）。\\
\begin{itemize}
  \item \textbackslash\{\}texttt\{train\textbackslash\{\}\_qwen35.py\}（QLoRA 4bit NF4 + ZeRO-2）就绪；gemma4-e4b 需 transformers 补丁复核。
  \item ssh5090 可访问后：\textbackslash\{\}texttt\{bash scripts/scp\textbackslash\{\}\_to\textbackslash\{\}\_ssh5090.sh\} + 数据同步 + 开训。
\end{itemize}

\subsection{8. verifiers 风格环境（P6）}

\textbackslash\{\}texttt\{src/experiments/game\textbackslash\{\}\_env.py\}：score\_trajectory 五轴 rubric（completion/progress/activity/\\
consistency/composite），可直接评分历史 run JSON，也可 GameEnv.rollout 在线跑分。\\
8 单测。为后续 RL 微调提供 env+rubric 抽象。\\

\subsection{9. 给同学四个问题的直接回答}

1. \textbackslash\{\}textbf\{空间一致性\}：版本化世界模型 + 失效检测（stale/面板/能力翻转）+ 分级恢复；\\
关键是「交互方式失效」要有一等公民的检测通道（本游戏是 failCount 翻转，其他游戏可能是\\
autoFishing 类能力标志）。\\
2. \textbackslash\{\}textbf\{时间一致性\}：三层时间戳对齐 + 阶段契约（precondition/allowed\_actions/success+settle/\\
timeout）+ 前缀 hash 防策略改写 + 失败分级。\\
3. \textbackslash\{\}textbf\{策略优化\}：引导式 prompt 的方向对，但先把\textbackslash\{\}textbf\{输出契约\}（token 预算、末行 JSON、\\
解析回退）做扎实；最优性评估用模拟器穷举做真值，别让模型自己评。\\
4. \textbackslash\{\}textbf\{效率\}：默认 L0 状态探针 + 触发器升级；DEBUG 环形内存、触发窗口落盘；\\
回退分级（行为/状态/策略）。\\

\subsection{10. 后续建议（优先级序）}

1. 面板处理泛化：把 LosePanel 个案抽象为「模态面板阻断」检测器（任意 Panel 激活+输入无效化），\\
进阶段契约守卫。\\
2. ssh5090 QLoRA：先 qwen3.5-4b（管线已验证），数据用 cold-start + processed-runs 合并。\\
3. 用 verifiers 环境批量产 failure\_recovery 数据（对抗性失败注入）。\\
4. 9B/E4B 基准补齐；VL struct 准确率上去后替代 mimo 视觉降本。\\
5. 22 游戏 benchmark：每游戏 profile + 3 seed × 2 配置（机制开/关）配对复跑。\\
\end{document}
